\section{임피던스의 역사적 기원}\label{sec:u3-history}

임피던스(impedance)라는 말은 어원상 진행을 방해한다는 뉘앙스를 갖지만,
공학에서는 단순한 방해를 넘어 시스템이 입력된 흐름에 어떻게 저항하고,
에너지를 저장했다가 되돌려주는지를 함께 나타내는 양이다.
이 개념이 왜 필요했는지는 19세기 후반 장거리 전신선과 전화선의 문제로
돌아가 보면 선명해진다.

\subsection{긴 선로가 깨뜨린 저항만의 그림}

책상 위의 짧은 직류 회로에서는 흐름을 방해하는 정도를 저항 하나에 담아
옴의 법칙 $V=IR$만으로 많은 것을 설명할 수 있다.
그러나 전보는 짧은 펄스가 켜졌다 꺼지는 신호이고, 전화는 사람 목소리처럼
계속 변하는 신호라서 직류가 아니다.
실제 긴 케이블에서는 이상한 일이 벌어졌다.
멀리 보낸 펄스는 늦게 도착하며 모양이 둥글게 퍼졌고, 목소리는 높은 성분과
낮은 성분이 서로 다르게 전달되어 멀리서 뭉개져 들렸다.
펄스의 날카로운 모서리일수록 높은 주파수 성분을 많이 필요로 하므로, 선로가
높은 성분을 더 약하게 만들거나 주파수마다 도착 시간을 다르게 만들면 점과
선의 구분이 흐려진다.
전선 저항만 줄이면 된다는 직류식 생각으로는 이 현상을 설명할 수 없었다.

이유는 긴 전선이 단순한 저항이 아니기 때문이다.
긴 선로에는 단위 길이당 저항 $R'~[\Omega/\mathrm{m}]$, 인덕턴스
$L'~[\mathrm{H/m}]$, 커패시턴스 $C'~[\mathrm{F/m}]$, 절연 누설을 나타내는
컨덕턴스 $G'~[\mathrm{S/m}]$가 위치를 따라 분포한다.
$R'$와 $G'$는 에너지가 열과 누설로 사라지는 손실 요소이고, $L'$와 $C'$는
자기장과 전기장에 에너지를 잠시 저장했다가 되돌려주는 저장 요소이다.
효과를 한 점에 모아 보는 집중 소자 모델과 달리, 긴 선로는 이런 작은
효과가 위치를 따라 반복해 붙어 누적되는 분포 시스템으로 보아야 한다.
위치 $x$와 시간 $t$에서의 전압 $v(x,t)$와 전류 $i(x,t)$로 이 생각을 적으면
전신방정식(telegrapher's equations)의 기본 모양이 된다.
\begin{align}
  \frac{\partial v}{\partial x} &= -R'i - L'\frac{\partial i}{\partial t}, \\
  \frac{\partial i}{\partial x} &= -G'v - C'\frac{\partial v}{\partial t}.
\end{align}
첫째 식은 선로를 따라 갈수록 전압이 변하는 이유가 직렬 저항 손실과 직렬
인덕턴스 전압에 있다는 뜻이고, 둘째 식은 전류가 변하는 이유가 병렬 누설
전류와 커패시터 충전 전류에 있다는 뜻이다.
지금 이 식을 풀 필요는 없다.
긴 전선은 한 덩어리 저항이 아니라 작은 저장소와 손실 요소가 끝없이 이어진
물체라는 표시로만 읽으면 충분하다.

\subsection{헤비사이드와 무왜곡 조건}

올리버 헤비사이드(Oliver Heaviside)는 전보 기사로 일하며 이 문제를 직접
접했고, 맥스웰의 전자기 이론을 통신선 문제에 적용하여 긴 선로의 왜곡을
수학적으로 다루었다 \cite{heaviside_electrical_papers,donaghyspargo2018heaviside}.
그의 통찰은 신호를 가장 세게 보내는 것이 목표가 아니라, 여러 주파수 성분이
서로 비슷한 규칙으로 약해지고 지연되게 만드는 것이 목표라는 데 있다.
파형은 여러 주파수 성분의 합이므로, 일부 성분만 더 약해지거나 더 늦어지면
전체 모양이 무너진다.
균일하고 선형이며 $R'$, $L'$, $C'$, $G'$가 주파수에 거의 의존하지 않는
이상화된 선로에서, 왜곡 없는 전달 조건은
\begin{equation}
  \frac{R'}{L'} = \frac{G'}{C'}
\end{equation}
로 주어진다.
$R'/L'$는 직렬 쪽에서 손실이 자기장 저장에 비해 얼마나 큰지,
$G'/C'$는 병렬 쪽에서 누설이 전기장 저장에 비해 얼마나 큰지를 나타내며,
두 양 모두 단위가 $\mathrm{s}^{-1}$이므로 이 조건은 같은 종류의 시간 척도를
맞추는 식이다.
오래된 전화선은 누설 $G'$가 작아 $G'/C'$가 작은데 인덕턴스가 부족해
$R'/L'$가 상대적으로 컸고, 일정 간격의 코일로 유효 $L'$를 키우면 두 비율이
음성 대역에서 가까워졌다.
전류 변화를 방해하는 인덕턴스를 오히려 추가해서 신호가 좋아지는 셈인데,
이는 방해를 키우는 일이 아니라 손실과 저장의 균형을 맞추어 파형 왜곡을
줄이는 설계이다.
그래서 헤비사이드가 1880년대에 사용한 impedance라는 말은 단순한 신조어가
아니라, 저항만으로 설명되지 않는 유도성과 용량성 효과까지 함께 보며 신호
전달을 설명하려는 실무 언어였다 \cite{ethw_heaviside,ptb_impedance_history}.

\subsection{저항, 리액턴스, 임피던스}

직류에서는 흐르기 어려운 정도를 저항 하나로 말해도 충분한 경우가 많다.
그러나 교류에서는 신호가 얼마나 작아지는지뿐 아니라, 전압과 전류의
흔들림이 시간상 얼마나 앞서거나 뒤지는지도 함께 보아야 한다.
전화선 한쪽에서 한 가지 주파수의 순음을 계속 보내고 반대쪽에서 얼마나
작아졌는지, 얼마나 늦게 흔들리는지를 재는 실용적 시험을 떠올리면 된다.
인덕터는 자기장에, 커패시터는 전기장에 에너지를 저장하며 주파수에 따라
전압과 전류 사이의 위상을 바꾸는데, 이런 저장 요소가 만드는 주파수
의존적인 방해를 리액턴스(reactance)라고 부른다.
리액턴스의 단위도 옴($\Omega$)이지만, 저항처럼 평균적으로 에너지를 열로
없앤다는 뜻은 아니다.
그래서 교류의 방해는 소산 성분과 저장 성분을 함께 담아
\begin{equation}
  Z = R + \jj X
\end{equation}
로 쓴다.
실수부 $R$은 에너지를 소산하는 저항이고, 허수부 $X$가 위상을 바꾸는
리액턴스이며, $\jj$는 $90^\circ$ 위상 차이를 표시하는 허수 단위이다.
정현파 정상상태에서 세 기본 소자의 임피던스는 각각
\begin{equation}
  Z_R = R, \qquad
  Z_L = \jj\omega L, \qquad
  Z_C = \frac{1}{\jj\omega C} = -\frac{\jj}{\omega C}
\end{equation}
이다($\omega=2\pi f$는 각주파수).
인덕터는 $\omega$가 커질수록 빠른 전류 변화에 더 큰 전압을 요구하고,
커패시터는 높은 주파수 전류를 상대적으로 잘 통과시킨다.
직렬 연결에서는 세 소자에 같은 전류가 흐르고 전원이 감당할 전압은 세 소자
전압의 합이므로, 전체 임피던스도 세 임피던스의 합이 된다.
\begin{equation}
  Z(\jj\omega) = Z_R + Z_L + Z_C
  = R + \jj\left(\omega L - \frac{1}{\omega C}\right).
\end{equation}
인덕터의 양의 리액턴스와 커패시터의 음의 리액턴스는 위상 방향이 반대라서
서로 빼지며, 어떤 주파수에서는 상쇄되어 회로가 저항만 남은 것처럼 보인다.
이 상쇄 현상이 공진의 씨앗인데, 공진 자체의 자세한 전개는 공개 확장판에서 다룬다.
이렇게 하면 미분방정식으로 풀어야 할 교류 회로도 $V=IZ$라는 옴의 법칙과
닮은 대수식으로 다룰 수 있게 된다.
전신선로 해석의 세부 유도와 페이저 표기의 확장 논의는 공개 확장판을
참고한다.

전보선 이야기가 남긴 질문은 시간에 따라 변하는 입력이 시스템을 지나며
크기와 위상이 어떻게 바뀌는가이며, 전기 회로와 기계 시스템을 같은 문장으로
비교하려면 먼저 이 질문을 적을 라플라스 변환과 선형 시불변 시스템의
문법이 필요하다.
